% !TEX root = ../main.tex
\documentclass[../main.tex]{subfiles}

\begin{document}

\intermezzo{Visual Materials: Mapping Embodied Identity}

\begin{center}
\begin{tikzpicture}
    \fill[AccentPurple] (-2,0) rectangle (-1.7,0.08);
    \fill[AccentPurple!70] (-1.5,0) rectangle (-1.2,0.08);
    \fill[AccentPurple!50] (-1,0) rectangle (-0.7,0.08);
    \fill[AccentPurple!30] (-0.5,0) rectangle (-0.2,0.08);
    \fill[AccentPurple!15] (0,0) rectangle (0.3,0.08);
\end{tikzpicture}
\end{center}

\vspace{0.5cm}

\begin{chapteroverview}
This chapter explores visual materials---logos, photography, video, and design elements---as sources for mapping \textbf{embodied identity}: how organizations materialize their roles through aesthetic choices. Images communicate identity dimensions that text cannot capture: professionalism, innovation, sustainability commitment, and institutional positioning.
\end{chapteroverview}


%═══════════════════════════════════════════════════════════════════════════════
\section*{Why Visual Analysis?}
%═══════════════════════════════════════════════════════════════════════════════

Visual elements carry meaning that complements textual analysis:

\sidenote{\textbf{Visual Signals:} Color palettes, logo design, photography style, infographic aesthetics, video production quality.}

\begin{description}[leftmargin=0.5cm, itemsep=0.3em]
\item[Immediate Impression] Visuals communicate before text is read
\item[Identity Clustering] Similar aesthetics signal similar positioning
\item[Professionalism Markers] Production quality signals resources and legitimacy
\item[Sector Alignment] Visual conventions reveal field membership
\end{description}


%═══════════════════════════════════════════════════════════════════════════════
\section*{Types of Visual Data}
%═══════════════════════════════════════════════════════════════════════════════

\subsection*{Logos and Branding}

Logos encode identity choices:
\begin{itemize}[leftmargin=*, itemsep=0.2em]
\item Color: Green for sustainability, blue for technology, warm tones for human focus
\item Typography: Modern sans-serif vs. traditional serif
\item Iconography: Abstract vs. representational
\item Complexity: Minimalist vs. detailed
\end{itemize}

\subsection*{Photography and Imagery}

Website and social media images reveal:
\begin{itemize}[leftmargin=*, itemsep=0.2em]
\item Team composition and diversity
\item Workspace environments
\item Technology and equipment
\item Partner and event documentation
\end{itemize}


%═══════════════════════════════════════════════════════════════════════════════
\section*{Analytical Approaches}
%═══════════════════════════════════════════════════════════════════════════════

\begin{description}[leftmargin=0.5cm, itemsep=0.3em]
\item[Computer Vision] Automated image classification, object detection, color analysis
\item[Embedding Similarity] Neural network embeddings for visual similarity clustering
\item[Qualitative Coding] Human interpretation of visual rhetoric and symbolism
\item[Multimodal Integration] Combining visual and textual signals
\end{description}


%═══════════════════════════════════════════════════════════════════════════════
\section*{Chapter Summary}
%═══════════════════════════════════════════════════════════════════════════════

\begin{summarybox}
\begin{itemize}[leftmargin=*, itemsep=0.3em]
\item \textbf{Visual materials} communicate identity dimensions that text cannot capture.

\item \textbf{Logos and branding} encode positioning through color, typography, and iconography.

\item \textbf{Photography} reveals team composition, resources, and institutional context.

\item \textbf{Analytical approaches} combine computer vision, embedding similarity, and qualitative interpretation.

\item \textbf{Integration} with textual modalities reveals consistency or divergence in identity presentation.
\end{itemize}
\end{summarybox}

\vspace{0.5cm}

\begin{center}
\begin{tikzpicture}
    \fill[AccentPurple] (-2,0) rectangle (-1.7,0.08);
    \fill[AccentPurple!70] (-1.5,0) rectangle (-1.2,0.08);
    \fill[AccentPurple!50] (-1,0) rectangle (-0.7,0.08);
    \fill[AccentPurple!30] (-0.5,0) rectangle (-0.2,0.08);
    \fill[AccentPurple!15] (0,0) rectangle (0.3,0.08);
\end{tikzpicture}
\end{center}

\closeintermezzo

\end{document}
